\renewcommand{\Anno}{2024}

\renewcommand{\NumeroProblema}{3}

\renewcommand{\Testo}{
La Scuola Sant’Anna deve assegnare due borse di studio per trascorrere un semestre all'estero. Il concorso, un po' sbrigativo, consiste nella risoluzione di un singolo esercizio, la cui risposta è un numero. Chi fornisce la risposta giusta guadagna un punto, chi fornisce la risposta sbagliata guadagna zero punti. Al termine, i due candidati con il punteggio maggiore vincono le due borse di studio. In caso di parità tra due o più candidati (sostanzialmente inevitabile visto che i punteggi finali possibili sono solo 0 e 1), il regolamento prevede che prevale sempre il candidato o i candidati più giovani.

Al concorso si presentano tre candidati, A, B, C, con A più giovane di B, e B più giovane di C.

\begin{enumerate}
\renewcommand{\labelenumi}{(\alph{enumi})}

\item  Supponiamo che, con la loro preparazione, tutti e tre i candidati abbiamo probabilità 1/3 di fornire la risposta giusta all'esercizio, indipendentemente l'uno dall'altro.

Determinare le probabilità di ciascuno dei tre candidati di vincere alla fine una delle due borse di studio, verificando in particolare che la loro somma è uguale a~2.

\item  Vistosi svantaggiato dall'età, il candidato C decide di non lavorare autonomamente all'esercizio, ma di copiare la risposta fornita da uno dei suoi due avversari.

Se copia da A, come si modifica la sua probabilità finale di ottenere una delle due borse di studio? E se invece copia da B?

(Ancora una volta si suppone che A e B abbiano probabilità 1/3 di fornire la risposta giusta, indipendentemente l'uno dall'altro)

\item  Supponiamo ora che i tre candidati abbiano tre probabilità $a$, $b$, $c$, eventualmente diverse tra di loro, di fornire la risposta corretta. Supponiamo inoltre che, per tali valori delle probabilità, i tre candidati abbiano la stessa probabilità di ottenere alla fine una delle due borse di studio.

Determinare i possibili valori di $b$ per cui questo può accadere.

\end{enumerate}
}

\renewcommand{\Risposta}{}

\renewcommand{\Aiutino}{}

\renewcommand{\SoluzioneUfficiale}{}

\renewcommand{\SoluzioniAlternative}{}

\renewcommand{\Osservazioni}{}

\renewcommand{\Commenti}{}
